\begin{conclusions}
El desarrollo de la presente investigación ha permitido materializar una plataforma integral para la 
gestión de imágenes dermatoscópicas dentro del ecosistema DermaUH, dando respuesta al problema científico 
planteado. En este sentido, se concluye que se cumplieron 
satisfactoriamente los objetivos generales y específicos trazados al inicio del proyecto.

El análisis exhaustivo del estado del arte de los principales repositorios internacionales de imágenes 
dermatológicas ---ISIC, HAM10000, BCN20000, PAD-UFES-20 y Derm7pt--- evidenció una desproporción 
sistemática a favor de fototipos claros (Fitzpatrick I--III), así como una concentración 
geográfica en centros de excelencia europeos, norteamericanos y australianos. Esta realidad fundamenta la 
necesidad imperativa de construir un repositorio soberano que capture la diversidad fenotípica de la 
población cubana, caracterizada por una prevalencia significativa de fototipos III, IV y V.

Para dar respuesta a esta necesidad, se diseñó y modeló un esquema de base de datos relacional que 
incorpora más de 30 atributos clínicos estructurados por registro, abarcando variables demográficas, 
antecedentes epidemiológicos, características clínicas de la lesión, clasificación diagnóstica con su 
método de confirmación e indicadores histológicos específicos para melanoma. Este esquema se alineó 
deliberadamente con la taxonomía del diccionario de datos ISIC v2.4, garantizando la interoperabilidad 
futura con repositorios internacionales y la generación de conjuntos de datos listos para el entrenamiento 
de modelos de inteligencia artificial.

La arquitectura del sistema se fundamentó en los principios de la Arquitectura Limpia, implementada 
mediante cuatro capas lógicas claramente diferenciadas ---Dominio, Aplicación, Infraestructura y 
Presentación--- que garantizan la separación de responsabilidades, la independencia tecnológica del 
núcleo de negocio y la escalabilidad del sistema ante futuras extensiones funcionales. Además, la 
plataforma integra un sistema de seguridad multicapa que combina autenticación mediante JSON Web Tokens, 
autenticación delegada mediante Google OAuth~2.0, control de acceso basado en 
roles y una política de seguridad por defecto (\textit{secure-by-default}) que minimiza la superficie de 
ataque y protege los datos clínicos sensibles.

Se implementó un motor de validación clínica con reglas condicionales que trasciende las restricciones 
convencionales de campos obligatorios, garantizando la coherencia médica de los metadatos mediante un 
árbol de decisión que refleja restricciones clínicas reales. Este mecanismo eleva la calidad del 
\textit{corpus} de datos y previene la persistencia de registros inconsistentes, aspecto crítico para la 
fiabilidad del entrenamiento de modelos de aprendizaje profundo. Complementariamente, se implementaron 
mecanismos de exportación masiva en formatos CSV y ZIP, acoplados a un flujo de autorización explícita, 
que habilitan la integración directa con los \textit{pipelines} de entrenamiento de modelos de 
inteligencia artificial, cumpliendo el objetivo estratégico de que los datos recopilados constituyan un 
activo de alto valor para la investigación computacional.

El módulo de estadísticas y analíticas desarrollado permite la generación de visualizaciones interactivas 
de distribuciones diagnósticas, demográficas y geográficas, proporcionando a los especialistas y 
administradores una herramienta de apoyo al diagnóstico y a la toma de decisiones epidemiológicas 
fundamentada en datos reales del territorio nacional.

La integración multimodal del sistema ---que articula una API REST documentada con el estándar OpenAPI, 
una interfaz web interactiva basada en Blazor Server y la capacidad de interoperar con aplicaciones 
móviles--- valida la viabilidad del enfoque de ciclo cerrado propuesto, donde la recolección de datos 
clínicos desde la consulta médica fluye hacia un repositorio estructurado y normalizado que retroalimenta 
la investigación y la capacitación.

En síntesis, la investigación ha demostrado que es factible diseñar e implementar una arquitectura 
flexible para una base de imágenes dermatoscópicas que integre metadatos clínicos locales y se acople 
fluidamente con las aplicaciones del ecosistema DermaUH, elevando la organización y trazabilidad de los 
datos dermatológicos en el país. El producto de \textit{software} resultante constituye un activo 
tecnológico soberano de incalculable valor para la Facultad de Matemática y Computación y los centros 
de salud asociados, sentando las bases para el desarrollo de una inteligencia artificial dermatológica 
verdaderamente equitativa y clínicamente pertinente en el contexto cubano.
\end{conclusions}
